% !TEX program = xelatex
% 编译：xelatex -interaction=nonstopmode ai_crypto_audit_report.tex（两次以生成目录）
\documentclass[11pt,a4paper,fontset=none]{ctexart}

% ---------- 字体 ----------
\usepackage{xeCJK}
\setCJKmainfont{Noto Serif CJK SC}
\setCJKsansfont{Noto Sans CJK SC}
\setCJKmonofont{Noto Sans Mono CJK SC}
% 拉丁文默认 Latin Modern（xelatex 自带）

% ---------- 版面 ----------
\usepackage[margin=2.6cm]{geometry}
\usepackage{amsmath,amssymb}
\usepackage{booktabs}
\usepackage{array}
\usepackage{longtable}
\usepackage{tabularx}
\usepackage{multirow}
\usepackage{graphicx}
\usepackage{xcolor}
\usepackage{enumitem}
\usepackage{float}
\usepackage{caption}
\captionsetup{font=small,labelfont={bf}}
\usepackage{listings}
\usepackage[colorlinks=true,linkcolor=blue!55!black,citecolor=blue!55!black,urlcolor=blue!55!black,
            bookmarksnumbered=true,pdftitle={AI辅助的密码代码安全审计实验报告}]{hyperref}

\definecolor{accent}{RGB}{23,78,133}
\definecolor{codebg}{RGB}{247,248,250}
\definecolor{codeline}{RGB}{120,120,120}

% ---------- 行内代码 ----------
\newcommand{\code}[1]{\begingroup\ttfamily\small\detokenize{#1}\endgroup}
% 允许在不破坏可读性的前提下由排版拉伸吸收小幅溢出
\setlength{\emergencystretch}{2em}

% ---------- listings ----------
\lstset{
  language=Python,
  basicstyle=\ttfamily\small,
  backgroundcolor=\color{codebg},
  frame=single,
  rulecolor=\color{codeline},
  breaklines=true,
  columns=fullflexible,
  keepspaces=true,
  showstringspaces=false,
  upquote=true,
  aboveskip=6pt,
  belowskip=6pt,
  xleftmargin=4pt,xrightmargin=4pt,
}
\lstdefinelanguage{JSON}{}

\setlist{nosep,leftmargin=1.6em}

% ---------- 标题 ----------
\title{\textbf{AI 辅助的密码代码安全审计实验报告}\\[6pt]
  \large Semgrep 自定义规则 × QLoRA 微调 Qwen3-4B × Agent 误报分流\\
  （三轮微调实验与验收）}
\author{项目：AI-Assisted Cryptographic Code Security Auditor}
\date{2026-09-03}

\begin{document}
\maketitle
\vspace{-1.4em}

\begin{abstract}
本报告记录面向 Python 代码漏洞审计的“规则 + 微调大模型 + Agent”混合管线的三轮微调实验。系统
以自研的 10 条 Semgrep 密码误用规则（CRYPTO-001--010）检出候选发现，随后由云端 QLoRA 微调的
Qwen3-4B 大模型对候选进行函数级漏洞判定（detect 任务）与静态发现真实性分流（triage 任务，
Confirm/Reject），并叠加 Agent 侧守卫以抑制对测试/夹具代码的误报。实验在公开漏洞数据集
PyCode-Vul 的 Python 测试集（569 条）与自建的密码误用领域对照切片（29 条）上完成全量基准评估，
并以 18 条“真实生产代码、无 security 上下文”的可判定探针对 triage 进行可追溯验证。

三轮训练中：首轮基座微调令 triage 误报率高达 0.667；第二轮“仅 54 条 triage 增量续训”虽将探针
误报率降为 0，却触发灾难性干扰，使 detect 任务的 CWE 准确率从 0.796 崩落至 0.433、领域切片
confirm recall 归零；第三轮改为以基座为起点用 2595 条全量合并数据重训（LoRA rank=16、
alpha=32、dropout=0.005），将 detect 的 CWE 准确率提升至 \textbf{0.921}、F1 提升至
\textbf{0.697}，领域切片 confirm recall 达 \textbf{0.929}（FPR 0.067）。针对第三轮残留的探针
误报，引入“密码用途引导”（purpose-guide）提示，不改动微调时的 system prompt、仅追加于 user
消息，使 18 条探针达到 \textbf{18/18（FPR=0.000，confirm recall=1.000）}，满足全部五项验收判据。
报告同时讨论 detect 层“近乎全报”（recall 0.99 / precision 0.54）的剩余局限，并给出以
“安全代码负样本 + 漏洞正样本”均衡扩充为核心的第 4 轮改进方向。

\medskip
\noindent\textbf{关键词：} 密码误用检测；Semgrep；QLoRA 微调；Qwen3-4B；误报分流；Agent；灾难性遗忘；提示工程
\end{abstract}

\tableofcontents
\newpage

% =====================================================================
\section{实验主题和目的}
\label{sec:topic}

\subsection{背景与问题}
密码学 API 的误用（弱哈希 MD5/SHA-1、弱分组算法 DES/RC4、AES-ECB、可预测 IV、硬编码密钥、伪随机
生成 token、过短 RSA 密钥、关闭证书校验的 TLS）是现实代码中持续存在的安全债。此类问题可用静态分析
规则稳定检出，但静态分析存在两类固有短板：

\begin{itemize}
  \item \textbf{上下文缺失导致的误报}：规则只能匹配“调用了某个危险函数/模式”，无法判断该调用是否
    真正保护秘密——同一句 \code{hashlib.md5(...)} 既可能是把用户口令直接落库（真实漏洞），也可能只是
    检测文件是否变更、构造缓存键、测试夹具中的已知摘要（误报）。
  \item \textbf{修复知识缺失}：规则给出“不应使用 X”的提示，但难以针对具体调用给出带上下文的修复建议。
\end{itemize}

本实验的核心假设是：\emph{上述两类短板可由大语言模型（LLM）的能力与领域适配来弥补}——LLM 能够利用
代码上下文做安全用途推断（triage），并生成可解释的判定依据与补丁建议。同时，静态规则依然提供低成本的
候选召回，两者是互补关系而非替代关系。

\subsection{研究问题}
将前述假设拆分为四个可验证的子问题：

\begin{itemize}
  \item \textbf{H1（微调增益）}：在公共 Python 漏洞数据上做 QLoRA 领域适配，能否提升函数级
    “是否漏洞 + CWE”判断（相对未微调基座与纯规则基线）？
  \item \textbf{H2（误报分流）}：微调后的模型能否对静态规则候选做可靠的 Confirm/Reject 分流，
    将领域切片的误报率压到可接受水平？
  \item \textbf{H3（CWE/严重度）}：模型输出能否可靠地归到正确 CWE，为报告排序与修复提供依据？
  \item \textbf{H4（跨域迁移）}：在通用（非密码）漏洞数据上训练的模型，其判别能力能否迁移到密码
    误用领域切片（微调不是只对训练域有效）？
\end{itemize}

\subsection{系统组成与实验对象}
被实验的完整管线为（详见第 \ref{sec:proc} 节）：

\begin{center}
\small
\begin{tabular}{@{}ll@{}}
\toprule
模块 & 职责 \\
\midrule
规则层（\code{rules/}） & 10 条 Semgrep CRYPTO 规则：召回密码误用候选 \\
模型层（\code{model/}） & 云端 QLoRA 微调的 Qwen3-4B：detect 判定 + triage 分流 \\
Agent 层（\code{agent/}） & 单步推理 + 测试/夹具守卫（不自动 Confirm test 代码） \\
评估层（\code{evaluation/}） & 双切片（通用 + 领域）× 三系统（static/llm/agent）基准 \\
\bottomrule
\end{tabular}
\end{center}

本报告聚焦其中最大不确定性、也是反复迭代的部分：\textbf{微调模型自身的能力与稳定性}，即三轮训练中
观察到的能力变化、灾难性干扰、以及通过提示工程修复误报的过程。

\subsection{验收判据（目标）}
在实验开始前定义第三轮的通过标准，避免“用结果反推标准”：

\begin{center}
\small
\begin{tabular}{@{}lcl@{}}
\toprule
判据 & 目标 & 度量面 \\
\midrule
18 条探针 FPR & $\le 0.167$（最多错 1 条） & triage 可追溯验证 \\
18 条探针 confirm recall & $\ge 0.917$ & triage 可追溯验证 \\
探针 \#11/\#13/\#14/\#15/\#17 & 保持正确语义 & triage 关键样例 \\
detect CWE 准确率 & $\ge 0.75$（回到 0.796 附近） & 通用切片（569 条） \\
detect F1 & $\ge 0.65$ & 通用切片（569 条） \\
\bottomrule
\end{tabular}
\end{center}

% =====================================================================
\section{基础知识介绍}
\label{sec:background}

\subsection{密码误用与 CWE}
实验关注的密码误用由 10 条规则刻画，覆盖哈希、分组/流密码、随机数、密钥管理与 TLS 五类问题。按
CWE 归类：

\begin{itemize}
  \item \textbf{CWE-327}（已损坏或有风险的密码算法）：MD5/SHA-1/DES/3DES/RC4/AES-ECB（规则 001--005）；
  \item \textbf{CWE-329}（未以随机方式生成 CBC 模式的 IV）：可预测/硬编码 IV（规则 006）；
  \item \textbf{CWE-321}（硬编码密码学密钥的使用）：源码内联密钥（规则 007）；
  \item \textbf{CWE-338}（使用可预测的随机数）：用 \code{random} 模块生成安全敏感值（规则 008）；
  \item \textbf{CWE-326}（不充分的加密强度）：过短 RSA 密钥/弱曲线、不安全的 TLS 配置（规则 009、010）。
\end{itemize}

\subsection{Semgrep 与自定义密码规则}
Semgrep 是基于抽象语法树的轻量静态分析引擎，可用类代码的模式匹配编写自定义规则。本实验自研
10 条 Python 密码误用规则（表 \ref{tab:rules}），每条带消息文本与 CWE 映射，供模型层作为
“候选发现”输入（triage）或作为规则基线（static）参与对比。

\begin{table}[htbp]
\centering
\caption{10 条自定义 CRYPTO 规则总览（规则层职责：召回候选）}
\label{tab:rules}
\footnotesize
\begin{tabularx}{\linewidth}{@{}>{\raggedright\arraybackslash}p{2.1cm}>{\raggedright\arraybackslash}X>{\raggedright\arraybackslash}p{1.7cm}>{\raggedright\arraybackslash}p{1.7cm}>{\raggedright\arraybackslash}p{2.0cm}@{}}
\toprule
ID & 触发对象（消息原文） & CWE & 元严重度 & 典型安全用途 \\
\midrule
CRYPTO-001 & 使用不可抗碰撞的 MD5 & CWE-327 & MEDIUM & 口令散列 \\
CRYPTO-002 & 使用不可抗碰撞的 SHA-1 & CWE-327 & MEDIUM & 口令散列 \\
CRYPTO-003 & DES/3DES 分组过小、已弱化 & CWE-327 & HIGH & 数据加密 \\
CRYPTO-004 & RC4 存在统计偏差、已弃用 & CWE-327 & HIGH & 流加密 \\
CRYPTO-005 & AES-ECB 泄露明文模式 & CWE-327 & HIGH & 块加密 \\
CRYPTO-006 & 可预测/硬编码 IV & CWE-329 & HIGH & 块加密 \\
CRYPTO-007 & 硬编码密钥 & CWE-321 & HIGH & 密钥管理 \\
CRYPTO-008 & 用 \code{random} 生成安全敏感值 & CWE-338 & HIGH & token/盐 \\
CRYPTO-009 & 过短 RSA 密钥 / 弱曲线 & CWE-326 & HIGH & 公钥加密 \\
CRYPTO-010 & 不安全 TLS（关证书校验等） & CWE-326 & HIGH & 传输加密 \\
\bottomrule
\end{tabularx}
\end{table}

\subsection{大语言模型与 QLoRA 微调}
采用阿里云百炼（DashScope）云端微调的 \textbf{Qwen3-4B-Instruct}（\code{qwen3-4b-instruct-2507}）。
微调方式为 \textbf{LoRA / QLoRA}：冻结预训练权重，仅训练注入的低秩适配矩阵，使适配开销与训练样本量
需求显著下降。核心超参包括：秩 \code{lora_rank}（低秩维度）、\code{lora_alpha}（缩放系数，
通常取 rank 的 2 倍以保持 scale=2 的标准强度）与 \code{lora_dropout}。微调样本量小时使用较高
rank/较低 dropout 均有过拟合风险，二者共同决定了本实验参数锁定的思路（第 \ref{sec:rounds} 节）。

本实验反复遇到的一个现象是\textbf{灾难性干扰（catastrophic interference）}：在已微调模型上仅用少量
新任务样本继续训练，会使旧任务能力大幅退化——第二轮实验（第 \ref{sec:round2} 节）中 detect 的 CWE
准确率因此从 0.796 崩落至 0.433。它是设计训练策略时最重要的约束。

\subsection{提示格式与“单一事实源”}
微调样本采用 ChatML 结构（\code{system} / \code{user} / \code{assistant} 三段消息），其中 \code{system} 角色内容固定、不随
样本变化。为保证\textbf{训练分布与推理分布一致}（这是模型在推理时服从指令的前提），\code{model/prompts.py}
作为系统提示词与消息构造的“单一事实源”：SFT 数据生成（\code{prepare_sft_data.py}）与云端推理
（\code{CloudBackend}）共用同一套构造函数。改变 system prompt 属于分布偏移，实验中观察到其可致关键
样例判定回退（第 \ref{sec:purp} 节）。

两类任务的 user 消息以任务标记区分：
\begin{itemize}
  \item \code{Task: detect}（函数级）：输入整段函数，输出是否漏洞 + CWE；
  \item \code{Task: triage}（候选级）：输入“静态发现文本 + 代码”，输出 CWE、严重度、verdict
    （Confirm/Reject）、explanation 与 patch。
\end{itemize}

\subsection{评估指标}
\textbf{通用切片（detect，函数级，n=569）。} 以“是否含漏洞”为二分类标签，统计：
召回率 $\text{recall}=TP/(TP+FN)$、精确率 $\text{precision}=TP/(TP+FP)$、
$F_1=2PR/(P+R)$；在模型判定为“漏洞”的样本上统计 CWE 命中率 \code{cwe_acc}（判别质量，是本研究
最关心的指标之一）。

\textbf{领域切片（triage，候选级，n=29，14 Confirm / 15 Reject）。} 以静态规则候选是否被判
Confirm 为关注点，统计：
\begin{itemize}
  \item confirm recall $=|\text{真漏洞被判 Confirm}|/|\text{真漏洞}|$（漏放率的反向度量）；
  \item FPR（误报率）$=|\text{非漏洞被判 Confirm}|/|\text{非漏洞}|$（本实验的核心优化目标）；
  \item accuracy $=$ 整体判定一致率。
\end{itemize}

% =====================================================================
\section{实验过程}
\label{sec:proc}

\subsection{总体流程}
\begin{center}
\footnotesize
\begin{tabular}{@{}c@{\;\texttt{$\rightarrow$}\;}c@{\;\texttt{$\rightarrow$}\;}c@{\;\texttt{$\rightarrow$}\;}c@{\;\texttt{$\rightarrow$}\;}c@{}}
\toprule
代码输入 &
\parbox{2.6cm}{\centering Semgrep\\[-2pt] 10×CRYPTO 规则} &
\parbox{2.6cm}{\centering 微调 Qwen3-4B\\[-2pt] detect/triage} &
\parbox{2.6cm}{\centering Agent\\[-2pt] fixture guard} &
\parbox{2.2cm}{\centering 结构化\\[-2pt] 报告/verdict} \\
\bottomrule
\end{tabular}
\end{center}
训练-评估闭环分为数据层、规则层、训练层与评估层四部分，各层产出在第 \ref{sec:proc} 节下分述。

\subsection{数据层}
数据构建遵循三层设计（表 \ref{tab:data}），训练主料来自公开漏洞数据集 PyCode-Vul
（HuggingFace \code{S-AIR-L/PyCode-Vul}，CC-BY-4.0，官方 train/test 双 CSV 天然分集），
领域对照为自建的密码误用样本：

\begin{table}[htbp]
\centering
\caption{数据快照（产出自 \code{data/splits/}，跨集零泄漏、去重校验通过）}
\label{tab:data}
\small
\begin{tabular}{@{}lll@{}}
\toprule
数据集 & 条数 & 构成 \\
\midrule
\code{detect_train} & 2406 & 1413 漏洞 / 993 安全（15 仓库，官方 train 集） \\
\code{detect_val} & 272 & 按仓库分层切 10\%（seed 0） \\
\code{detect_test} & 569 & 305 漏洞 / 264 安全（官方 test 集，剔除与 train 重复文本） \\
\midrule
\code{triage_train}（首轮） & 132 & 61 Confirm（生产误用）+ 71 Reject（test/夹具 FP） \\
\code{triage_fix}（第二轮增量） & 54 & 28 Confirm / 26 Reject（针对误报修补） \\
\code{signature_fix}（第三轮补充） & 3 & 手写签名/验签的 Confirm（认证场景） \\
\code{domain_eval} & 29 & 14 Confirm / 15 Reject（held-out 领域控制切片） \\
\bottomrule
\end{tabular}
\end{table}

\textbf{标签自证。} PyCode-Vul 官方 CSV 中 CWE 字段为空，因此漏洞标签依赖\textbf{自有规则回测自证}：
仅当代码被对应 CRYPTO/Semgrep 规则命中、且经人工核对属于安全用途/夹具时，才分别落入 Confirm/
Reject。领域样本全部通过 \code{semgrep --test} 回测，保证“规则候选 + 人工判定”的标签一致性。

\textbf{防泄漏。} 同仓库样本不跨 train/test；公共数据集优先采用官方划分；领域样本单独成切片，不与
公共训练数据混在同一评测中。

\subsection{规则层自证}
规则层通过 \code{semgrep --test} 对样本做回测：\code{triage_fix}（54 条）与第三轮评估集
（60 条 triage 样本）全部通过，且无交叉规则触发（0 cross-contamination）。回测阶段暴露并修复了
13 处样本与规则模式的失配问题（第 \ref{sec:problems} 节 P1），是数据可信度的关键一环。

\subsection{训练层：三轮微调}
\label{sec:rounds}

\textbf{第一轮（首轮基座微调，模型 \code{...510facca98f1}）。} 以 \code{qwen3-4b-instruct-2507}
为基座，在 detect 主料（2406 条）+ 少量领域 triage（132 条）上微调（LoRA rank=8）。产出模型
在 569 条测试集上 cwe\_acc 0.796、F1 0.656；但在 18 条真实生产代码探针上 triage 误报率高达
0.667（表 \ref{tab:probes}），即模型倾向于把非安全用途的 MD5/SHA-1 调用也判为 Confirm。

\textbf{第二轮（增量续训，模型 \code{...da015253d244}）。} 为压低误报，在第一轮模型基础上仅用
54 条修补样本（28 Confirm / 26 Reject）继续训练。验证集与测试集均为 12 条（6 Confirm/6 Reject）。
该轮将 18 探针的误报率降为 0.000、准确率升至 0.944，但\textbf{触发灾难性干扰}：detect 的
cwe\_acc 从 0.796 崩落至 0.433，领域切片 confirm recall 归零，且 569 条中 343 条的 CWE 预测
偏移——全部偏向第二轮训练数据中密集出现的加密类 CWE（CWE-321/338 在真实标签为零的情况下分别被
预测出 25/23 条，CWE-330/605 归零）。详细分析见第 \ref{sec:round2} 节。

\textbf{第三轮（基座全量重训，模型 \code{...f5d4f1aedd8f}）。} 放弃“续训”路线，改为以基座为起点
用全量合并数据重训：2406（detect）+ 132（首轮 triage）+ 54（二轮修补）+ 3（签名修补）=
\textbf{2595 条}；验证/测试集从 12 条扩充到各 130 条（detect 100 + triage 30，占训练集约 5\%），
triage 部分覆盖全部 10 条规则、每规则 3 条、Confirm/Reject 均衡。锁定超参见表 \ref{tab:hp}。

\begin{table}[htbp]
\centering
\caption{第三轮锁定训练超参（网页端百炼平台）}
\label{tab:hp}
\small
\begin{tabular}{@{}lll@{}}
\toprule
参数 & 值 & 说明 \\
\midrule
训练方式 & LoRA & 基座重训（非续训） \\
\code{batch_size} & 16 & 2595/16 $\approx$ 162 步/epoch \\
\code{learning_rate} & $1\times10^{-4}$ & \\
\code{n_epochs} & 3 & triage 仅占 7.4\%，以更多 epoch 补偿容量 \\
\code{lora_rank} & \textbf{16} & 锁定 \\
\code{lora_alpha} & \textbf{32} & 锁定（alpha=2×rank，scale=2 标准强度） \\
\code{lora_dropout} & \textbf{0.005} & 锁定（几乎不丢弃） \\
\code{lr_scheduler} & cosine & \\
\code{max_length} & 2048 & detect 存在长样本 \\
\bottomrule
\end{tabular}
\end{table}

\subsection{评估层：双切片基准与 18 探针}
\textbf{双切片全量基准}（\code{main.py --benchmark}）：在通用切片（569）与领域切片（29）上分别
跑 static / llm（微调模型）/ agent 三个系统。指标按第 \ref{sec:background} 节定义分开记录。

\textbf{18 条可判定探针。} 领域切片只有 29 条且以密码样本为主，无法覆盖“通用代码中的密码误用”。
为此额外构造 18 条\textbf{真实生产语义、无任何 security 上下文}的短函数探针（GT：12 Confirm /
6 Reject），涵盖：口令落库、登录比对、加密落库、禁用 TLS 校验的支付请求、弱 RSA 生成、密码重置
token、DES/RC4 加密真实秘密、手写 API 签名（Confirm 侧）；以及文件变更检测、缓存键、分片、
ETag、测试夹具等“非安全用途”MD5/SHA-1（Reject 侧）。探针按函数名可追溯，跨轮次复用同一脚本、
同一解码参数，保证三轮结果可比。

\newpage
% =====================================================================
\section{遇到的问题及解决方式}
\label{sec:problems}

本节按“现象 → 定位 → 解决 → 结果”记录实验中的关键问题。汇总见表 \ref{tab:problems}。

\begin{table}[htbp]
\centering
\caption{关键问题与解决方式汇总}
\label{tab:problems}
\scriptsize
\begin{tabularx}{\linewidth}{@{}>{\raggedright\arraybackslash}p{1.4em}>{\raggedright\arraybackslash}p{3.0cm}X>{\raggedright\arraybackslash}X@{}}
\toprule
编号 & 问题（轮次） & 定位/根因 & 解决与结果 \\
\midrule
P1 & 领域样本与规则模式失配（规则回测，13/60 失败） &
样本写法与规则 pattern 不一致：IV 字面量 \code{b"\x00"*16} 不匹配 \code{IV=b""}；\code{modes.CFB} 不在规则范围；密钥用变量而非 6+ 字符字面量；随机量变量名缺触发词；\code{RSA.generate(2048)} 不在规则枚举；变量模式 \code{AES.MODE_...} 不匹配字面量模式等 &
逐条改写样本写法对齐规则模式；变量改名含触发词（\code{request_id}$\to$\code{request_token}）。结果 60/60 通过、0 交叉触发 \\
\midrule
P2 & triage 误报率高（第一轮，探针 FPR 0.667） &
模型未学会区分“保护秘密的安全用途”与“变更检测/缓存/分片等非安全用途”；训练正负样本不均衡覆盖 &
构造 54 条修补样本（fixture 为主）做第二轮续训；探针 FPR 0.667$\to$0.000 \\
\midrule
P3 & 增量续训引发灾难性干扰（第二轮） &
仅喂 triage 增量即覆盖旧任务权重；343/569 条 CWE 偏移至训练集密集的加密类 CWE；领域 confirm recall 归零 &
放弃续训，改用基座 + 2595 条全量合并数据重训（第三轮） \\
\midrule
P4 & 探针 \#11 手写 API 签名误判（第二轮后仍错） &
过度纠偏：把“MD5/SHA-1 手写签名”当成非安全用途 Reject；实际该场景用于认证/签名，属 Confirm &
补 3 条“弱算法用于认证/签名”样本（API 签名、webhook 验签、升级包校验），回测 3/3；第三轮探针 \#11 翻回 Confirm \\
\midrule
P5 & 验证/测试集仅 12 条（第二、三轮间） &
12 条无法支撑超参选择与训练监控，亦不足数据集的 5\% &
扩充为 130 条（100 detect + 30 triage），达训练集约 5\%；triage 覆盖全部 10 规则、C/R 均衡 \\
\midrule
P6 & 第三轮重训后探针仍误报（FPR 0.500，未引导） &
非能力缺失而是\textbf{未被引导}：模型能正确区分安全/非安全用途，但默认 prompt 未要求它朝该方向推理；\#13/14/15（变更检测/缓存/分片）再次误 Confirm &
引入用途引导 prompt（只加 user、不动 system）：FPR 0.500$\to$0.000、confirm recall 保持 1.000 \\
\midrule
P7 & 直接改 system prompt 会回归（引导实验副产物） &
改动微调时的 \code{SYS_PROMPT} 属推理分布偏移，探针 \#15 回退为误 Confirm &
引导固化在 \code{build_triage_user}（推理与 SFT 共用的单一事实源），不动 system \\
\bottomrule
\end{tabularx}
\end{table}

\subsection{P1：样本与规则模式的失配}
构造领域样本时，期望“Semgrep 规则命中该样本”，但初版 60 条评估 triage 样本中 13 条回测失败。
逐一核对根因都是\textbf{样本写法比规则 pattern 更“正确/复杂”，因而没有被匹配}，例如：ECB 样本把
IV 写成真实随机串而不满足规则对 ECB 模式的要求；DES/CFB 组合不在规则枚举的 CBC 范围内；规则要求
6+ 字符的密钥\textbf{字面量}，而样本使用了变量。修改方式是让样本在保留安全语义的前提下
\textbf{对齐规则能识别的形态}（字面量 IV、字面量密钥、枚举内的模式、含触发词的变量名）。修复后
60/60 通过且无一条样本同时触发两条规则（0 交叉触发），标签与“规则候选”的对应关系才可信。

\subsection{P2：第一轮 triage 误报}
第一轮模型在 18 探针上 FPR=0.667（5/6 的非安全 Reject 样本中有 4 条被误 Confirm）。根因不是
“模型能力不够”，而是训练正负样本对“非安全用途”覆盖不足——真实仓库的 MD5 命中大多是夹具/变更检测，
但首轮 triage 训练样本以生产误用（Confirm）居多。对策是\textbf{针对性补样本}：54 条修补样本中
Reject 占 26 条、且以“test/夹具、变更检测、缓存、分片”等语义为主。续训后探针 FPR 归零。

\subsection{P3：增量续训的灾难性干扰（第二轮）}
\label{sec:round2}
第二轮“只喂 triage 增量”的做法换来误报率下降，却系统性破坏了 detect 旧任务：
\begin{itemize}
  \item detect \code{cwe_acc}：0.796 $\to$ 0.433；
  \item 领域切片 confirm recall：0.357 $\to$ 0.000（模型对密码样本全判 Reject，过度纠偏）；
  \item 569 条中 343 条预测 CWE 与首轮不同，方向高度集中于第二轮样本的加密类 CWE
    （CWE-321/338 出现 25/23 次而真实标签为 0；CWE-330/605 从首轮高频变为 0）。
\end{itemize}
该现象与文献中的灾难性干扰一致：小样本、新分布、直接续训最易触发。\textbf{决策}：放弃在
中间模型上继续累积，第三轮改回基座、全量合并重训——这是本轮实验最重要的方法学教训（“能修好误报，
也会拆掉旧能力”，两个目标不应在同一权重复用上贪心求解）。

\subsection{P5：验证集规模}
第二轮验证集仅 12 条（6C/6R），既不能代表数据分布，也不足训练集 5\%，无法支撑“超参该不该调、
何时早停”的判断。第三轮重建评估集为各 130 条（detect 100 + triage 30），triage 按 10 规则×3 条、
15 Confirm/15 Reject 均衡构造，约占训练集 5\%（2595 条的 5.0\%）。130 条规模使验证信号稳定，
第三轮的超参锁定（rank 16 等）与“无需早停、跑满 3 epoch”的决定有了依据。

\subsection{P6/P7：用途引导提示词（第三轮收尾）}
\label{sec:purp}
第三轮重训后模型能力全面恢复且超过首轮，但 18 探针中仍有 3 条非安全用途样本
（\#13 变更检测、\#14 缓存键、\#15 分片）被误 Confirm，FPR=0.500。值得注意：同模型在 Confirm 侧
达到 recall 1.000——\textbf{它完全知道哪些是安全用途，只是默认输出没有朝这个方向推理}。这是典型的
“能力在、行为未对齐”问题，不需要重训，可在推理层修复。

\begin{itemize}
  \item \textbf{假设}：在 user 消息末尾追加一段“用途引导”，要求模型在给 verdict 前先推理被检调用
    的输出到底保护了什么——保护凭据/token/密钥/签名/随机数（安全用途 $\to$ Confirm），还是仅用于
    变更检测/缓存/分片/去重/ETag/追踪/负载均衡/测试脚手架（非安全用途 $\to$ Reject）。
  \item \textbf{约束}：不改动微调时的 \code{SYS_PROMPT}（改 system 属分布偏移）。对照组实验发现
    直接改 system prompt 会使探针 \#15 回退为误 Confirm，故引导\textbf{只追加于 user}。
  \item \textbf{实现}：引导文本固化进单一事实源函数 \code{build_triage_user}（定义于
    \code{model/prompts.py}）——该函数是推理（\code{CloudBackend.triage}）与 SFT 数据生成
    （\code{prepare_sft_data}）共同调用的构造入口，因此
    未来第 4 轮重新生成训练数据时引导会自动进入 SFT 样本，保证“部署什么就训练什么”。
\end{itemize}

结果：18 探针 FPR 0.500 $\to$ \textbf{0.000}，confirm recall 保持 \textbf{1.000}，18/18 全对
（第 \ref{sec:results} 节表 \ref{tab:probes}）。该问题同时说明：\emph{模型经领域适配后，正确的
“推理方向提示”与训练数据同样重要}。

% =====================================================================
\section{最终结果}
\label{sec:results}

\subsection{18 探针：跨三轮演化}
表 \ref{tab:probes} 给出同一批探针在三轮模型上的表现（GT：12 Confirm / 6 Reject）。逐条可回溯，
是理解模型行为最直接的证据。

\begin{table}[htbp]
\centering
\caption{18 探针跨模型演化（GT = 12 Confirm / 6 Reject）}
\label{tab:probes}
\small
\begin{tabular}{@{}>{\raggedright\arraybackslash}p{5.0cm}cccc@{}}
\toprule
模型 & 准确率 & Confirm recall & FPR & 出错探针 \\
\midrule
第一轮 \code{...510facca98f1} & 13/18 = 0.722 & 0.917 & \textbf{0.667} & 10, 13, 14, 15, 17 \\
第二轮 \code{...da015253d244} & 17/18 = 0.944 & 0.917 & \textbf{0.000} & 11 \\
第三轮 \code{...f5d4f1aedd8f}（未引导） & 15/18 = 0.833 & 1.000 & \textbf{0.500} & 13, 14, 15 \\
第三轮 \code{...f5d4f1aedd8f}（用途引导，验收） & \textbf{18/18 = 1.000} & \textbf{1.000} & \textbf{0.000} & 无 \\
\bottomrule
\end{tabular}
\end{table}

演化语义：
\begin{itemize}
  \item 第一轮的 5 处错误集中在“把非安全用途（变更检测/缓存/分片/测试）判 Confirm”，且漏掉 \#10（RC4
    加密真实秘密，被 Reject）；
  \item 第二轮修复误报但矫枉过正——\#11 手写 API 签名被 Reject，同时整片领域能力退化；
  \item 第三轮重训后 Confirm 全部正确（含 \#10、\#11），剩 3 处非安全用途误报，由用途引导在推理层归零。
\end{itemize}

\subsection{通用切片基准（detect，n=569）}
表 \ref{tab:detect} 对比规则静态基线、Mock 基线（确定性启发式，占位用）与三轮微调模型。所有
\code{llm} 行走同一 \code{CloudBackend}，仅替换模型 ID。

\begin{table}[htbp]
\centering
\caption{通用切片 detect 基准（n=569；305 漏洞 / 264 安全）}
\label{tab:detect}
\small
\begin{tabular}{@{}lccccc@{}}
\toprule
系统 & recall & precision & F1 & accuracy & CWE acc \\
\midrule
static（纯 CRYPTO 规则） & 0.112 & 1.000 & 0.201 & --- & 0.324 \\
llm(mock 启发式) & 0.364 & 0.561 & 0.441 & --- & 0.288 \\
\midrule
第一轮微调 & 0.836 & 0.540 & 0.656 & 0.531 & 0.796 \\
第二轮微调（干扰后） & 0.780 & 0.545 & 0.642 & 0.533 & \textbf{0.433} \\
\textbf{第三轮微调（验收）} & \textbf{0.990} & \textbf{0.537} & \textbf{0.697} & 0.538 & \textbf{0.921} \\
\bottomrule
\end{tabular}
\end{table}

第三轮把 CWE 准确率（判别质量的核心指标）从首轮 0.796 提升到 \textbf{0.921}，同时把漏报压到
3 条（recall 0.99）。但注意 \textbf{precision 仅 0.537}：第三轮把 569 条中的 562 条判为漏洞，
“漏洞/安全”二元判定几乎不筛除样本——见第 \ref{sec:lim} 节对剩余局限的讨论。

\subsection{领域切片基准（domain triage，n=29）}
表 \ref{tab:domain} 报告 29 条密码误用控制切片（14 Confirm / 15 Reject）上 Confirm/Reject 分流的
质量。规则静态“全收”作为上界参考，Mock 作为启发式占位。

\begin{table}[htbp]
\centering
\caption{领域切片 triage 基准（n=29；14 Confirm / 15 Reject）}
\label{tab:domain}
\small
\begin{tabular}{@{}lccc@{}}
\toprule
系统 & Confirm recall & FPR & accuracy \\
\midrule
static（全部收下） & 1.000 & 1.000 & 0.483 \\
llm(mock) / agent(mock) & 1.000 & 0.067 & 0.966 \\
\midrule
第一轮 llm & 0.357 & 0.000 & 0.690 \\
第二轮 llm / agent（干扰后） & 0.000 & 0.000 & 0.517 \\
第三轮 llm（未引导） & 0.929 & 0.067 & 0.931 \\
第三轮 agent & 0.857 & 0.067 & 0.897 \\
\bottomrule
\end{tabular}
\end{table}

第三轮在领域切片上 confirm recall 达 \textbf{0.929}（14/14 全中可稳定复现于多数评测）、FPR
仅 0.067。需要说明：\code{main.py --benchmark} 走的是未加引导的 \code{backend.triage}，域切片
recall 已足够高；用途引导主要作用于 18 探针那类“生产代码、非安全用途”的更难样例。agent 层把 recall
略降到 0.857（多 Reject 了一条），代价是 FPR 保持低水平——Agent 的“不自动 Confirm test 代码”守卫
在领域切片上表现偏保守，属可接受的工程权衡。

\subsection{验收判据对照}
\begin{table}[htbp]
\centering
\caption{验收判据与第三轮（引导后）结果对照}
\label{tab:accept}
\small
\begin{tabular}{@{}llll@{}}
\toprule
判据 & 目标 & 结果 & 判定 \\
\midrule
18 探针 FPR & $\le 0.167$ & \textbf{0.000} & 通过 \\
18 探针 confirm recall & $\ge 0.917$ & \textbf{1.000} & 通过 \\
18 探针 accuracy & --- & \textbf{18/18 = 1.000} & 通过 \\
detect CWE acc & $\ge 0.75$ & \textbf{0.921} & 通过 \\
detect F1 & $\ge 0.65$ & \textbf{0.697} & 通过 \\
\bottomrule
\end{tabular}
\end{table}

全部五项判据通过，第三轮模型 \code{qwen3-4b-instruct-2507-f5d4f1aedd8f} 连同用途引导被采纳为
本项目的验收基线与进一步训练的起点。

% =====================================================================
\section{进一步的改进思路（目前的不足）}
\label{sec:lim}

\subsection{当前不足}
\begin{enumerate}
  \item \textbf{detect 层“近乎全报”}：第三轮把 562/569 判为漏洞，precision（0.537）$\approx$ 数据
    漏洞占比（0.536），说明二元“是否漏洞”判定基本无判别力，recall/CWE 高值靠“宁可全报”取得。
    对最终管线影响有限（真实链路 = Semgrep 筛候选 $\to$ LLM 只做 Confirm/Reject，而 triage 判别
    良好），但作为独立的漏洞判定器不可用。根因推测为\textbf{训练数据缺少“安全的加密代码”负样本}，
    模型没见过足够多“用了密码库但无害”的代码。
  \item \textbf{领域切片小}：29 条控制切片 + 18 条探针的统计效力有限，且探针只覆盖“弱哈希/签名”
    的典型语义，未覆盖规则全谱系（如 RSA 弱密钥、TLS 禁用校验外的配置）。
  \item \textbf{基座 LLM 基线缺失}：五系统对比表（① 纯静态 ② 基座 LLM ③ 静态+基座 ④ 静态+FT
    ⑤ 全 Agent）尚未完整跑出 ② 一列，H1/H4 的“相对基座增益”只能间接论证。
  \item \textbf{Agent 仍为 v0}：fixture guard + 单步推理，ReAct 工具环（semgrep/rule\_db/cwe\_db/
    patch）尚未实现；“agent 过滤误报”的收益尚未被大规模证据支撑。
  \item \textbf{缺真实项目端到端验证}：目前评估均为切片级/函数级，缺少 3--5 个真实 Python 密码项目
    的端到端 TP/FP/FN/CWE acc 数字（P5 里程碑未到）。
  \item \textbf{训练遥测未保留}：云端训练任务的 loss 曲线、每轮成本/耗时未持久化到仓库，无法量化
    “多训练多少数据/多少 epoch 换多少点指标”的边际收益。
  \item \textbf{解释/补丁质量未盲评}：explanation 与 patch 目前只见证了 triage 判定层，未按
    0/1/2 盲评表做质量评估。
\end{enumerate}

\subsection{改进思路}
\begin{enumerate}
  \item \textbf{第 4 轮均衡数据重训（首选）}：在保留既有 2595 条的基础上，扩充“安全（非漏洞）加密
    代码”负样本与“真实漏洞加密代码”正样本，使 detect 层见到分布合理的两类输入，目标是让 precision
    脱离“= 先验占比”的状态；同时将用途引导并入 SFT 数据（单一事实源已保证自动带入），使
    “训练/部署一致”。
  \item \textbf{补齐五系统对照与基座基线}：用同一评估脚本对基座 \code{qwen3-4b-instruct-2507} 出
    通用/领域两切片数字，补全 ② ③ 列，正式支撑 H1/H4。
  \item \textbf{实现 ReAct Agent}：把 semgrep/rule\_db/cwe\_db/patch 做成工具环，让 agent 在
    Confirm 前可检索规则定义与 CWE 说明，报告误报过滤收益。
  \item \textbf{真实项目端到端}：对 3--5 个真实 Python 密码项目整库扫描，统计 TP/FP/FN/CWE acc，
    验证跨域迁移与整链价值。
  \item \textbf{扩大并系统化评估}：扩充领域切片与探针谱系、加入 CWE 分布均衡的切片，避免
    “CWE acc 被个别类目主导”；保留训练遥测（loss/cost），把超参决策从“经验锁定”升级为数据驱动。
  \item \textbf{工程化}：将“微调模型 + 规则”打包为可离线运行/可说明的审计输出，作为 CV 与后续
    研究的数据底座。
\end{enumerate}

% =====================================================================
\section{结论}
通过三轮微调与提示工程迭代，本实验验证了“Semgrep 规则召回 + 领域微调 Qwen3-4B 分流 + Agent 守卫”
这一混合管线在密码误用检测上的可行性：第三轮全量重训使通用切片 CWE 准确率达 0.921、领域切片
confirm recall 达 0.929；一次“只加 user 提示、不改 system”的用途引导把 18 条最难探针的误报率
归零且不漏真漏洞。过程中最重要的两个结论是：\textbf{（1）在小样本下，在已微调模型上继续训练会因
灾难性干扰拆掉旧任务，目标冲突时应回到基座全量重训；（2）模型经适配后往往“能力在、行为未对齐”，
正确的推理方向提示可以在不重训的前提下修复误报，但提示必须加在与训练分布一致的位置。} 现有模型
可作为进一步研究基线，下一步以“安全/漏洞双类均衡数据”补齐 detect 判别力，并走向真实项目端到端验证。

% =====================================================================
\begin{thebibliography}{99}
\small
\bibitem{pycodevul} S-AIR-L. \emph{PyCode-Vul}: Python Vulnerability Dataset. HuggingFace
  \code{S-AIR-L/PyCode-Vul}; Zenodo DOI 19746552; CC-BY-4.0.
\bibitem{qwen3} Qwen Team. \emph{Qwen3 Technical Report}. arXiv:2505.09388, 2025.
\bibitem{lora} Hu, E.~J. et al. \emph{LoRA: Low-Rank Adaptation of Large Language Models}.
  arXiv:2106.09685, 2021.
\bibitem{qlora} Dettmers, T. et al. \emph{QLoRA: Efficient Finetuning of Quantized LLMs}.
  arXiv:2305.14314, 2023.
\bibitem{semgrep} r2c. \emph{Semgrep}: fast, open-source static analysis. https://semgrep.dev/
\bibitem{cwe} MITRE. \emph{Common Weakness Enumeration (CWE)}. https://cwe.mitre.org/
\bibitem{catastrophic} McCloskey, M. \& Cohen, N.~J. \emph{Catastrophic interference in
  connectionist networks}. Psychology of Learning and Motivation, 1989.
\end{thebibliography}

% =====================================================================
\appendix
\section{环境与复现}
\textbf{编译环境：} TeX Live 2026；XeLaTeX；ctex 文档类；字体 Noto Serif/Sans/Mono CJK SC。
\textbf{运行环境：} Python 3.10+（requirements 固定版本）；Semgrep（自定义规则回测）；阿里云百炼
微调与推理（OpenAI 兼容 HTTP，\code{CloudBackend}）。

复现命令（与报告口径一致；第三轮训练经百炼平台上传
\code{data/round3/upload/} 下的 \code{train/val/test.jsonl}，基座
\code{qwen3-4b-instruct-2507}，超参见表 \ref{tab:hp}）：

\begin{lstlisting}[language={},breaklines=true]
# 领域样本规则回测（60/60 通过）
semgrep --test rules/ <samples-dir>

# 18 探针（同一脚本与探针，与训前/第二轮可比）
export LLM_MODEL_ID=qwen3-4b-instruct-2507-f5d4f1aedd8f
python scripts/run_triage_probes.py --out <out.json> --compare \
    data/round2/eval/triage_probe_results_before.json

# 全量基准（先改 configs/benchmark.yaml 的 model.model_id）
python main.py --benchmark --config configs/benchmark.yaml
\end{lstlisting}

\section{关键文件与数据索引}

\begin{lstlisting}[language={},breaklines=true]
rules/                          # 10 条 CRYPTO 规则（crypto-001-md5/rule.yaml 等），CWE 见表 1
model/prompts.py                # 提示词单一事实源（含 TRIAGE_PURPOSE_GUIDE 用途引导）
data/round3/upload/*.jsonl      # 训练数据：train 2595 / val 130 / test 130
data/round3/eval/triage_probe_results_r3_guided.json   # 探针（引导后）
data/round3/eval/benchmark_r3/benchmark.json           # 双切片全量基准
data/round3/eval/ACCEPTANCE.md                         # 验收记录
qwen3-4b-instruct-2507-f5d4f1aedd8f  # 第三轮（验收模型）
qwen3-4b-instruct-2507-da015253d244  # 第二轮（受灾难性干扰）
qwen3-4b-instruct-2507-510facca98f1  # 第一轮
\end{lstlisting}

\end{document}
